\chapter{Conclusion}
\label{chap:conclusion}

% Sources: outputs/*/README.md, docs/Design Principles.md,
% docs/Justification of Implementation.md §2

\section{Summary}
\label{sec:conc:summary}

\glsreset{dag}
\glsreset{hrgs}
\glsreset{rgs}

This thesis formalized the schedule space left informal by the source papers~\cite{Benchasattabuse2026Bridging,Benchasattabuse2025Integrating} and searched it systematically, addressing the four contributions set out in Section~\ref{sec:intro:contributions}.

The \Gls{dag}-based model of Chapter~\ref{chap:model} provides a common framework for describing and evaluating purification schedules. Its evaluator reproduces the source paper's fidelity results to four significant figures (Section~\ref{sec:exp:validation}), allowing different purification patterns, copy counts, and heralding placements to be evaluated under the same physical model. Built on this model, the three-tier search framework of Chapter~\ref{chap:algorithms} finds schedules with higher rates than the paper's hand-picked example, both at matched cost and, more substantially, when the budget is relaxed (Section~\ref{sec:exp:headline}). Since the hand-picked schedule is presented by its authors as a demonstration rather than an optimized design, this comparison primarily validates that the search produces valid, working schedules rather than constituting the thesis's strongest result.

The stronger result is instead a feasibility advantage. Under inner-qubit error, which the source papers do not model, the paper's fixed schedule can fall below $F\ge0.9$ entirely, while a searched schedule at the same cost remains feasible (Section~\ref{sec:exp:netgen}). The same pattern appears across 100 random per-hop heterogeneous networks and in a deliberately constructed weak-link case (Section~\ref{sec:exp:heterogeneous}). These results show why schedule optimization becomes important when different parts of a chain experience different conditions: a fixed construction cannot redistribute purification effort, whereas a searched schedule can.

Chapter~\ref{chap:discussion} combines this pattern with the noise and chain-length sweeps to identify a common regime: the advantage is largest when the fidelity constraint is strict enough to make the choice of schedule matter, but not so strict that all searched schedules become infeasible. This regime appears across all three tested axes: noise level, chain length, and network heterogeneity. The timing experiments provide a separate observation: the preferred heralding strategy depends on the balance between communication delay, memory dephasing, and purification success probability, rather than being intrinsic to the purification operation itself.

\paragraph{Results Existential Claims Limitations}

These results are existence claims within the searched schedule families, not proofs of global optimality over the full legal schedule space (Section~\ref{sec:disc:scope}). The finite beam width and restricted construction families \emph{prevent concluding that no better schedule exists}. The excluded-construction check of Appendix~\ref{app:results} illustrates this directly: it identifies a valid, feasible schedule that no search tier can construct because its construction lies outside their search spaces, regardless of how good the schedules found within those spaces are.

Taken together, the results support the working hypothesis stated in Section~\ref{sec:intro:gap}: systematic search can improve on the paper's hand-picked baseline at equal cost and, more importantly, retain feasibility where that baseline fails. The practical implication is that treating purification scheduling as an optimization problem, rather than fixing it by convention, can reveal useful constructions that a fixed uniform design would miss.

\section{Future Work}
\label{sec:conc:future}

The results also point to several extensions that would broaden the scope of the optimizer and reduce the gap between the current model and a deployable network.

\begin{itemize}

      \item \emph{Adaptive scheduling.} Conditioning later purification decisions on earlier measurement or heralding outcomes, within a single execution, could adapt to the realized state of the network rather than committing to a fixed sequence in advance. As noted in Section~\ref{sec:model:scope}, this would most naturally be formulated as an \Gls{mdp}~\cite{Puterman2009Markov,Sutton2014Reinforcement}, extending the current non-adaptive schedule space and potentially avoiding purification steps that become unnecessary after an intermediate outcome. Reinforcement learning is one candidate solution method, though the state and action spaces would need to be redefined carefully for this setting.

      \item \emph{Online optimization.} Rather than adapting within one execution, the search itself could be re-run periodically as noise and link quality drift over minutes to months, instead of being fixed once at commissioning as it is throughout this thesis. This would allow the selected schedule to follow changes in physical conditions, particularly in heterogeneous networks where different hops may degrade at different rates.

      \item \emph{Multi-path routing.} Extending the \Gls{dag} formalism from a single repeater chain to tree or mesh networks would require a more general connectivity representation and a corresponding extension of the current span-based search. In such networks, purification scheduling and route selection would become coupled decisions: the optimizer could potentially choose not only where and when to purify, but also which physical path or set of paths should be used to establish the end-to-end pair.

      \item \emph{Beyond the Pauli-channel model.} Hardware-specific noise models, including more detailed photon-loss and gate-error processes, would let the optimizer be evaluated against a broader range of physical implementations and could change the relative ranking of candidate schedules.

\end{itemize}

\vspace{0.5cm}

Overall, this work provides a starting point for systematic purification-schedule optimization. Even within a restricted non-adaptive space, purification placement and allocation affect both rate and the ability to meet a target fidelity. Extending the approach to adaptive decisions, time-varying networks, richer topologies, and more realistic noise models would test whether these gains persist under more realistic conditions.

\vspace{1cm}

\paragraph{Code and data availability}
\label{sec:exp:availability}

The evaluator, the three search tiers of Chapter~\ref{chap:algorithms}, and every experiment script behind the tables and figures in this chapter and its appendix are openly available at
\url{https://github.com/EliottFlechtner/entanglement-purification-scheduler}.